\section*{Solutions to Homework 9}

\textbf{Part A.}

\begin{lstlisting}
$ bash codeHints/L09/setup_hw9.sh
$ cd practice
$ git clone hw9_server.git hw9
$ cd hw9
$ bash test.sh
FAIL: got 'Helo, world'
$ sed -i.bak 's/Helo/Hello/' greet.sh && rm greet.sh.bak    # or edit by hand
$ bash test.sh
PASS
$ git commit -am "Fix greeting typo" -m "Closes #1"
$ mkdir -p .github/workflows .github/ISSUE_TEMPLATE
\end{lstlisting}

Create \code{.github/workflows/ci.yml} (spaces only):

\begin{lstlisting}
name: CI
on:
  push:
    branches: [main]
  pull_request:
jobs:
  test:
    runs-on: ubuntu-latest
    steps:
      - uses: actions/checkout@v4
      - name: Run the tests
        run: bash test.sh
\end{lstlisting}

Create \code{.github/ISSUE\_TEMPLATE/bug\_report.md} (the front matter with \code{name:} and \code{about:} and the two headings of the lecture) and \code{.github/pull\_request\_template.md} (the heading and the checklist of the lecture). Then:

\begin{lstlisting}
$ git add .github
$ git commit -m "Add CI workflow and templates"
$ git tag v0.1.0 $(git rev-list --max-parents=0 HEAD)
$ git tag -a v1.0.0 -m "First release"
$ git push origin main v0.1.0 v1.0.0
$ printf 'next-bugfix=1.0.1\nissue-closes-when=merged\n' > ../hw9_answers.txt
\end{lstlisting}

(Committing the workflow and the templates in two commits, as the task text says, works equally well.)

\textbf{9.3--9.4.} The checker runs your project's \code{test.sh} on a copy of the server's \code{main}, which is exactly what the workflow does on GitHub. \code{git commit -m ... -m ...} makes the second \code{-m} a separate paragraph of the message; GitHub looks for the keyword and the issue number anywhere in it.

\textbf{9.5.} Without \code{uses: actions/checkout} the virtual machine would have no files, so \code{bash test.sh} would fail. The workflow file must be in \code{.github/workflows} (with an \emph{s}).

\textbf{9.7.} The lightweight tag is a bare name for the first commit; the annotated tag has its own object with your message, so \code{git show v1.0.0} shows a \emph{Tagger}. Tags are only pushed when you name them (or use \code{--tags}).

\textbf{9.8.} Fixing a bug in \code{1.0.0} gives the patch version \code{1.0.1}. \code{Closes \#12} closes the issue when the change is \emph{merged} into the default branch, not when the pull request is opened.

\textbf{Part B.}

\begin{lstlisting}
(GitHub) New repository: ci-practice, Public, no README
$ git remote add github git@github.com:YOURNAME/ci-practice.git
$ git push github main v0.1.0 v1.0.0
(GitHub) Actions tab: the "CI" workflow ran: green tick
$ printf 'repo=YOURNAME/ci-practice\n' >> ../hw9_answers.txt
(GitHub) Issues > New issue "Add a note to the README"  -> number 1
$ git switch -c readme-note
$ echo "Notes coming soon." >> README.md
$ git commit -am "Add a note to the README"
$ git push github readme-note
(GitHub) Open the pull request, description "Closes #1"; wait for the green check; Merge
(GitHub) Issue #1 is now closed
$ printf 'issue=1\n' >> ../hw9_answers.txt
(GitHub) Releases > Draft a new release > tag v1.0.0 > Generate release notes > Publish
\end{lstlisting}

Replace \code{YOURNAME} with your GitHub user name. The checker needs the answers file to contain \code{repo=}, \code{issue=}, and a remote of \code{practice/hw9} that points at that repository. It then asks GitHub's public interface whether the issue is closed, whether a release exists for \code{v1.0.0}, and whether the latest Actions run on \code{main} succeeded. If the last check is red, open the run's log in the Actions tab.

\Needspace{8\baselineskip}
\textbf{Common mistakes.}
\begin{itemize}
  \item A \emph{tab} character in the workflow (the checker reports it) or a folder called \code{.github/workflow} instead of \code{workflows}.
  \item Creating Issue \#1 \emph{before} the first push: the old commit message \code{Closes \#1} then closes it at once, and the pull request closes nothing. The check passes anyway, but you missed the exercise. Push first, then create the issue.
  \item Forgetting to push the tags: \code{git push origin main} sends no tags.
  \item Creating the GitHub repository with a README: the first push is rejected because the histories differ. Create it empty.
  \item A private repository: GitHub's public interface cannot see it, so the checker reports that it cannot find it.
  \item Writing \code{issue-closes-when=opened}: the issue closes when the change is \emph{merged}.
\end{itemize}
